\section{Pointers toward an alternative empirical signature}
\label{sec:pointers}

The present paper makes a deliberately narrow negative claim. A
negative result without a positive complement is not, on its own,
a scientific contribution; we therefore sketch in this section
the alternative empirical signature we develop in companion
work. The sketch is short because the development is the subject
of a separate methodological paper currently in preparation; we
refer the reader to that paper for the calibration, theoretical
scaffolding and falsifiable predictions, and limit ourselves
here to indicating what it contains.

\paragraph{A five-family non-cyclical signature.}
On the same panels and at the same significance level, five
diagnostic families consistently reject the AR(1) null: long
memory (Hurst exponent estimated by detrended fluctuation
analysis \citep{peng1994,hurst1951}), multifractality (MF-DFA
spectrum width \citep{kantelhardt2002,bacry_muzy_delour2001}),
non-linearity (BDS independence test
\citep{brock_dechert_scheinkman1996}), information complexity
(permutation entropy \citep{bandt_pompe2002}) and reflexive
drift (sliding-window Kolmogorov-Smirnov tests on
regime-conditional distributions). The joint signature
\(C+B+D+I+S\) rejects on a majority of cells on which the
cycle-based protocol rejects.

\paragraph{Generative models.}
Multifractal random walks \citep{bacry_muzy_delour2001,
calvet_fisher2002,calvet_fisher2008} provide a generative model
that reproduces the empirical signature without invoking any
canonical periodicity. Adaptive markets \citep{lo2017} and the
free-energy principle \citep{friston2010} provide a behavioural
rationale for the reflexive component. Random matrix theory
\citep{marchenko_pastur1967,laloux1999,bouchaud_potters2003}
characterises the cross-sectional structure.

\paragraph{Status of the companion paper.}
At the time of writing the companion paper is in preparation. We
acknowledge that the present negative result gains scientific
value once the alternative is on record, and we will not seek
journal submission of the present paper until the companion is
available as a citeable preprint. The two papers are therefore
intended to be read in conjunction.
